% Part:sets-functions-relations
% Chapter: size-of-sets
% Section: pairing-alt

\documentclass[../../../include/open-logic-section]{subfiles}

\begin{document}

\olfileid{sfr}{siz}{pai-alt}

\olsection{વૈકલ્પિક જોડ-નિરૂપણ વિધેય}

\begin{explain}
$\Nat^2$ની બીજી પરિગણનાઓ પણ છે જેમનાં પ્રતિવિધેયો
શોધવાં વધુ સરળ છે. અહીં એક આપીએ.
પરિગણનાને સરણિમાં જોવાને બદલે, શરૂઆતમાં ખાલી
જગ્યાઓ સાથે સંકળાયેલી ધન પૂર્ણાંકોની યાદી લો.
આ જગ્યાઓને $\tuple{n,m}$ જોડો વડે નીચે મુજબ
ક્રમશઃ ભરવાની કલ્પના કરો. પહેલા સ્થાને~$0$ હોય
એવી જોડોથી (એટલે કે $\tuple{0,m}$ જોડોથી) શરૂ કરો.
પ્રથમ જોડ (એટલે કે $\tuple{0,0}$) પ્રથમ ખાલી
જગ્યામાં મૂકો, પછી એક ખાલી જગ્યા છોડો, બીજી જોડ
(એટલે કે $\tuple{0,2}$) પછીની ખાલી જગ્યામાં મૂકો,
ફરી એક જગ્યા છોડો, અને એમ આગળ વધો.
હવે પરિગણનાની (અધૂરી) શરૂઆત આ રીતે દેખાય છે:
\[\small
\begin{array}{@{}c c c c c c c c c c c@{}}
\mathbf 1 & \mathbf 2 & \mathbf 3 & \mathbf 4 & \mathbf 5 & \mathbf 6 & \mathbf 7 & \mathbf 8 & \mathbf 9 & \mathbf{10} & \dots \\ \\
\tuple{0,0} &  & \tuple{0,1} &  & \tuple{0,2} &  & \tuple{0,3} & & \tuple{0,4} &  & \dots \\
\end{array}
\]
હજી ખાલી રહેલી જગ્યાઓ માટે $\tuple{1,m}$ જોડો
સાથે આ પ્રક્રિયા ફરી કરો, ફરીથી એકાંતરે ખાલી
જગ્યા છોડતા જાઓ:
\[\small
\begin{array}{@{}c c c c c c c c c c c@{}}
\mathbf 1 & \mathbf 2 & \mathbf 3 & \mathbf 4 & \mathbf 5 & \mathbf 6 & \mathbf 7 & \mathbf 8 & \mathbf 9 & \mathbf{10} & \dots \\ \\
\tuple{0,0} & \tuple{1,0} & \tuple{0,1} &  & \tuple{0,2} & \tuple{1,1} &
\tuple{0,3} & & \tuple{0,4} &  \tuple{1,2} & \dots \\
\end{array}
\]
એ જ રીતે $\tuple{2,m}$, $\tuple{3,m}$ વગેરે
જોડો મૂકો. આમ, પૂર્ણ કરેલી પરિગણનાની શરૂઆત
આ રીતે થાય છે:
\[\small
\begin{array}{@{}cc c c c c c c c c c@{}}
\mathbf 1 & \mathbf 2 & \mathbf 3 & \mathbf 4 & \mathbf 5 & \mathbf 6 & \mathbf 7 & \mathbf 8 & \mathbf 9 & \mathbf{10} & \dots \\ \\
\tuple{0,0} & \tuple{1,0} & \tuple{0,1} & \tuple{2,0}  & \tuple{0,2} &
\tuple{1,1} & \tuple{0,3} & \tuple{3,0}  & \tuple{0,4} &  \tuple{1,2} & \dots \\
\end{array}
\]
\sourcecorrection{OLSIZ-003}{મૂળની \texttt{pairing-alt.tex}, પંક્તિઓ 39--45માં
\(\tuple{2,m}\) બે વાર લખાયું છે. તરત પછીના કોષ્ટકમાં
\(\tuple{2,0}\) અને \(\tuple{3,0}\) આવતા હોવાથી અહીં બીજી જોડને
\(\tuple{3,m}\) કરી છે.}
ઉપરની સરણિના ખાનાંને આ પરિગણના પ્રમાણે અંકો
આપીએ, તો સુઘડ આડીઅવળી રેખાને બદલે આ ગોઠવણી મળે:
\[
\begin{array}{ c | c | c | c | c | c | c | c }
& \mathbf 0 & \mathbf 1 & \mathbf 2 & \mathbf 3 & \mathbf 4 & \mathbf 5 & \dots \\
\hline
\mathbf 0 & 1 & 3 & 5 & 7 & 9 & 11 & \dots \\
\hline
\mathbf 1 & 2 & 6 & 10 & 14 & 18 & \dots & \dots \\
\hline
\mathbf 2 & 4 & 12 & 20 & 28 & \dots & \dots & \dots \\
\hline
\mathbf 3 & 8 & 24 & 40 & \dots & \dots & \dots & \dots \\
\hline
\mathbf 4 & 16 & 48 & \dots & \dots & \dots & \dots & \dots \\
\hline
\mathbf 5 & 32 & \dots & \dots & \dots & \dots & \dots & \dots \\
\hline
\vdots & \vdots & \vdots & \vdots & \vdots & \vdots & \vdots & \ddots\\
\end{array}
\]
જોઈ શકાય છે કે હાર~$0$ની જોડો પરિગણનાના અયુગ્મ
અંકવાળાં સ્થાનો પર છે: જોડ $\tuple{0,m}$નું સ્થાન
$2m+1$ છે. બીજી હારની જોડો $\tuple{1,m}$ એવા
સ્થાનો પર છે જેમનો અંક અયુગ્મ સંખ્યાનો બમણો છે,
ખાસ કરીને $2 \cdot (2m+1)$.
ત્રીજી હારની જોડો $\tuple{2,m}$ એવા સ્થાનો પર
છે જેમનો અંક અયુગ્મ સંખ્યાનો ચાર ગણો છે,
$4 \cdot (2m+1)$; અને એમ આગળ વધે છે.
દરેક હાર માટે $(2m+1)$ના ગુણકો
$1$, $2$, $4$, $8$, \dots એ ચોક્કસપણે
$2$ના ઘાત છે: $1= 2^0$, $2 = 2^1$,
$4 = 2^2$, $8 = 2^3$, \dots\@
હકીકતમાં, સંબંધિત ઘાતાંક હંમેશાં તે જોડનો પ્રથમ
સભ્ય છે. તેથી જોડ $\tuple{n,m}$ માટે ગુણક
$2^n$ છે. આથી સામાન્ય સૂત્ર $2^n \cdot (2m+1)$
મળે છે. જોકે, આ જોડોનું \emph{ધન} પૂર્ણાંકોમાં
નિરૂપણ છે: $\tuple{0,0}$નું સ્થાન~$1$ છે.
સ્થાન~$0$થી શરૂ કરવું હોય, તો પરિણામમાંથી~$1$
બાદ કરવો પડે. આથી મળે છે:
\end{explain}

\begin{ex}
નીચે મુજબ આપેલું વિધેય $h\colon \Nat^2 \to \Nat$:
\[
h(n,m) = 2^n (2m+1) - 1
\]
પ્રાકૃતિક સંખ્યાઓની જોડોના ગણ~$\Nat^2$ માટેનું
જોડ-નિરૂપણ વિધેય છે.
\end{ex}

\begin{explain}
આ પ્રમાણે, $\Nat^2$ની આપણી બીજી પરિગણનામાં
જોડ $\tuple{0,0}$ની સંકેતસંખ્યા
$h(0,0) = 2^0(2\cdot 0+1) - 1 = 0$ છે;
$\tuple{1,2}$ની સંકેતસંખ્યા
$2^{1} \cdot (2 \cdot 2 + 1) - 1 = 2
\cdot 5 - 1 = 9$ છે; અને $\tuple{2,6}$ની
સંકેતસંખ્યા $2^{2} \cdot (2
\cdot 6 + 1) - 1 = 51$ છે.
\end{explain}

ક્યારેક પ્રાકૃતિક સંખ્યાઓની જોડો~$\Nat^2$ને
સંકેતબદ્ધ કરવું પૂરતું હોય છે; સંકેતન વ્યાપ્ત હોય
તે જરૂરી નથી. આવાં સંકેતનોનાં પ્રતિવિધેયો
માત્ર આંશિક વિધેયો હોય છે.

\begin{ex}
નીચે મુજબ આપેલું વિધેય $j\colon \Nat^2 \to \Nat^+$:
\[
j(n,m) = 2^n3^m
\]
એ !!a{injective} વિધેય $\Nat^2 \to \Nat$ છે.
\end{ex}

\end{document}
